\chapter{Tracing source rows and pixel addresses}
\MARSChapterHeading{Source rows and pixel addresses}
\label{revision:lab-arrays}

Use \texttt{sprite-rows.asm} from \texttt{revision/labs/files} with this chapter.
It pairs the main book's byte representation, row renderer and compact-branch
caller. These activities use normal MARS Run and Debug; the sixth revised unit
is not yet a native interactive lesson.

Your task is to connect a source bit, a destination coordinate and an actual word
store. Keep the row's base address separate from the renderer's moving argument
pointer. A drawing that looks plausible can still use the wrong stride, skip
black stores or read a neighboring object as another source row.

\section{Predict one asymmetric row}

Write \texttt{0x16} as eight bits, including the leading zeros. List its set
columns under the convention that bit seven is leftmost. Reverse the pattern and
list the mirrored row's set columns. Add eight to those local column numbers to
place them in the right half of the framebuffer.

Derive the addresses for the first pixel of each half and the first pixel of the
next row. Explain why the two displacements are 32 and 64 bytes. Then derive the
last pixel's starting displacement and the last byte of that store. A word-sized
access has a beginning and an end.

\section{Pause before the first renderer call}

Break at \texttt{row\_reversed}. The next instruction stores the computed reversed
byte; that store has not executed yet. Inspect the original row in
\texttt{\$s4}, the result in \texttt{\$v0}, the reversed-byte destination in
\texttt{\$s1}, and the framebuffer row base in \texttt{\$s2}. Predict which
objects remain unchanged before the pending byte store.

Continue to \texttt{draw\_row}. On the first visit, the arguments describe the
right half: pattern \texttt{0x68} and destination \texttt{frame}+32. Step to
\texttt{store\_pixel}, before its first word store. The first tested source bit
is zero, so the selected color is black. Predict the next destination address
and count after that store and the following updates.

At \texttt{right\_row\_drawn}, the helper has returned. Its argument pointer
has advanced to \texttt{frame}+64; the saved row base still points to
\texttt{frame}. The next call prepares a new argument for the left half. Explain
why the advanced argument alone is not the next iteration's complete state.

Continue to \texttt{row\_drawn}, before the outer updates. Both halves of row zero
are complete, while the source and output-row pointers still name row zero.
After the updates, they name row one. Check the remaining count at the next visit
to \texttt{row\_loop}. State the execution boundary when recording a pointer or
count; the same register can legitimately describe different stages of an iteration.

Figure~\ref{fig:revision-sprite-progress} connects the first completed row with
the finished pair. Compare these pictures with your recorded pointers and count;
the picture alone does not show whether the next outer update has executed.

\begin{figure}[htbp]
\centering
\begin{tikzpicture}[alt={Two native Bitmap Display observations. At row_drawn the first row contains the original and mirrored bit patterns and all lower rows are still black. At frame_complete all eight source rows are drawn beside their horizontal reflections. A publication divider marks the boundary between columns seven and eight.}]
\node[anchor=north west,inner sep=0,font=\small\bfseries] (firsttitle) at (0,0) {1. First row complete, before the outer updates};
\node[anchor=north west,inner sep=0] (first) at ([yshift=-2mm]firsttitle.south west) {\MARSRevisionCrop{revision/captures/sprite/sprite-first-row.png}{0.86\linewidth}{0.032031250}{0.737777778}{0.032031250}{0.178888889}};
\node[anchor=north west,inner sep=0,font=\small\bfseries] (fulltitle) at ([yshift=-7mm]first.south west) {2. All eight rows complete};
\node[anchor=north west,inner sep=0] (full) at ([yshift=-2mm]fulltitle.south west) {\MARSRevisionCrop{revision/captures/sprite/sprite-complete.png}{0.86\linewidth}{0.032031250}{0.155555556}{0.032031250}{0.178888889}};
\begin{scope}[shift={(full.south west)},x={($(full.south east)-(full.south west)$)},y={($(full.north west)-(full.south west)$)}]
\draw[black,line width=1.8pt] (0.5,0) -- (0.5,1);
\draw[white,line width=0.6pt,dashed] (0.5,0) -- (0.5,1);
\end{scope}
\end{tikzpicture}
\par\smallskip
{\small Left: original, columns 0--7.\quad Right: mirror, columns 8--15.}
\par\medskip
\begin{minipage}{\linewidth}
\small
\begin{enumerate}
\item At \texttt{row\_drawn}, both halves of row zero are drawn. The next instruction advances the source pointer. The strip enlarges that row; the lower seven rows in this fresh session remain black.
\item At \texttt{frame\_complete}, the pending instruction selects the exit service. The pixel stores are complete. The divider and labels are native publication annotations, not additional pixels written by the program.
\end{enumerate}
\end{minipage}
\caption{Two automated light-theme Debug captures of the worked program. Every logical pixel was checked against independently derived RGB values at each stop.}
\label{fig:revision-sprite-progress}
\end{figure}


\section{Separate row order from bit order}

Replace the source bytes with a single-bit pattern: row zero contains
\texttt{0x80}, row one \texttt{0x40}, continuing down to \texttt{0x01} in row
seven. Predict the diagonal in the original image and its horizontal reflection
in the right image. Unlike a symmetric pattern, this input reveals which bit and
which row supplies each visible position.

Next use an all-zero row followed by an all-one row. Confirm that the first row
is black even if its framebuffer words held nonzero values before drawing.
Skipping a black store would behave like transparency and violate this renderer's
contract. A new session's initial zeros can conceal that mistake.

A horizontal mirror reverses the bits within each byte. A vertical flip reverses
the order in which the bytes supply rows. Combining the two performs both changes.
Do not diagnose a vertically misplaced pattern by changing its within-row mask.

\section{Investigate an address error}

In a copy, change the outer row-base increment from 64 to 32. Predict which row's
words the next iteration will overwrite. Run a short Debug trace and locate the
first destination span that differs from your address calculation. The source
bytes and reversed bytes may still be correct because their pointers use separate
updates.

Restore the stride, then change a pixel increment from four bytes to one. This
moves through overlapping byte ranges rather than consecutive words. Release 6
ordinary misaligned accesses need not stop the program with an alignment fault;
inspect the actual bytes and colors instead of relying on an exception to detect
the mistake.

For another copy, omit the initialization of black before testing a bit. Follow
the color register's most recent writer. A white pixel can then leave white in the
register for a later clear bit. The first wrong color is a state-lifetime error,
not a different definition of the source pattern.

\section{Draw a vertical variation independently}

The separate \texttt{vertical-sprite.asm} caller supplies a source of eight bytes
and a sixteen-by-eight framebuffer. Implement \texttt{draw\_vertical}, replacing
the failing \texttt{break} body. Its arguments are the source address in
\texttt{\$a0} and the frame base in \texttt{\$a1}. It draws eight rows into
columns eight through fifteen, with source row seven supplying output row zero
and source row zero supplying output row seven. Within each row, bit seven remains
leftmost.

Preserve columns zero through seven, both guards and all source bytes. Restore
the saved registers and stack pointer from the chapter's calling convention,
return through the incoming continuation, and perform no console or device I/O
other than the ordinary framebuffer writes. The supplied \texttt{draw\_row}
helper is available; another equivalent renderer is allowed.

Plan the source row address, destination row base, count and return address before
writing instructions. If you call the helper, decide which state must survive
that call. A backwards source cursor may become one byte before the source after
the final update; your stopping condition must prevent another load there.

Run the public checks using JDK~25:

\begin{verbatim}
java -jar /path/to/cli.jar grade \
  revision/labs/checks/vertical-sprite.kts \
  /path/to/your/vertical-sprite.asm
\end{verbatim}

The checks replace the source and prefill the frame with varied values. They
compare all destination words, preserve the left half, source and guards, and
check saved registers and the restored stack pointer. They do not prescribe your
scratch registers, private labels or frame layout. Keep your own address derivation
and a trace of the first and final rows alongside a passing report.

The later clipped-blitter lab adds partial source and destination intersections.
Before attempting it, be able to explain this fully in-bounds operation for an
arbitrary source pattern and a prefilled destination. Clipping changes which
rows and columns are admitted; it does not change the size of a pixel word or
make an out-of-range address safe to access.
